\section{Related Work}
\label{sec:related-work}
\label{app:related-work} % FIX: intro references \ref{app:related-work}

\parbench{} draws on three lines of work: code translation benchmarks, LLM-for-HPC evaluation, and robustness testing.

% \paragraph{General-purpose code benchmarks.}
% Execution-based evaluation is now the standard for measuring functional correctness, but existing task distributions do not cover parallel HPC translation. HumanEval~\cite{HumanEval2021} and TransCoder~\cite{TransCoder2020} operate at function granularity, while SWE-bench~\cite{SWEbench2024} evaluates repository-level issue resolution in general software projects. These benchmarks do not exercise GPU memory models, thread synchronization, API-specific launch structure, or host-device coordination. SWE-bench Illusion~\cite{SWEbenchIllusion2025} further demonstrates that static benchmark scores can overstate model capability when evaluation methodology is not carefully controlled, a concern that motivates \parbench{}'s conjunctive build-run-verify pipeline and augmentation robustness probes.

\paragraph{General-purpose code benchmarks.}
Execution-based evaluation is the standard for measuring functional correctness, but existing benchmarks do not cover parallel HPC translation. HumanEval~\cite{HumanEval2021} and TransCoder~\cite{TransCoder2020} operate at function granularity, while SWE-bench~\cite{SWEbench2024} targets repository-level issue resolution—neither exercises GPU memory models, thread synchronization, or host-device coordination. SWE-bench Illusion~\cite{SWEbenchIllusion2025} further shows that static scores can overstate capability when evaluation methodology is not carefully controlled, motivating \parbench{}'s conjunctive build-run-verify pipeline and augmentation robustness probes.

% \paragraph{Parallel-code benchmarks.}
% ParEval~\cite{ParEval2024} and PCEBench~\cite{chen2025pcebench} study synthesis of parallel code from natural-language descriptions, testing whether a model can produce a correct solution but not whether it can preserve an existing implementation's decomposition across APIs. ParEval-Repo~\cite{ParEvalRepo2025} and RepoTransBench~\cite{Wang2026RepoTransBench} study full-repository translation and demonstrate a severe integration bottleneck: models fail once they must recreate build systems and repository scaffolding. AlphaTrans~\cite{AlphaTrans2025} introduces compositional neuro-symbolic translation at repository scale but does not target parallel APIs. \parbench{} is complementary: it preserves executable evaluation but fixes build infrastructure so that the measured object is kernel translation. VibeCodeHPC~\cite{VibeCodeHPC2025} evaluates prompt-driven HPC code generation, while QiMeng-MuPa~\cite{QiMengMuPa2025} addresses multi-paradigm parallel code synthesis; both focus on generation rather than translation of existing parallel implementations.

\paragraph{Parallel-code benchmarks.}
ParEval~\cite{ParEval2024} and PCEBench~\cite{chen2025pcebench} evaluate synthesis from natural-language descriptions, testing correctness but not cross-API preservation of existing implementations. ParEval-Repo~\cite{ParEvalRepo2025} and RepoTransBench~\cite{Wang2026RepoTransBench} study full-repository translation and reveal a severe integration bottleneck at build systems and scaffolding. AlphaTrans~\cite{AlphaTrans2025} introduces compositional neuro-symbolic translation at repository scale but does not target parallel APIs. VibeCodeHPC~\cite{VibeCodeHPC2025} and QiMeng-MuPa~\cite{QiMengMuPa2025} focus on HPC code generation rather than translation of existing implementations. \parbench{} is complementary to all of the above: it fixes build infrastructure so that the measured object is kernel translation fidelity.

% \paragraph{HPC-specific LLM evaluation.}
% LASSI~\cite{LASSI2024} is the nearest prior work in granularity and verification. It evaluates an agentic self-correction pipeline on 10 HeCBench kernels across two directions (CUDA to OpenMP-target and reverse), measuring how iterative repair improves translation success. \parbench{} asks a different question: what is the single-attempt translation capability across a broad set of directions, without iterative repair? It also differs in scope---35~kernels from five suites, six standard translation directions among CUDA, OpenMP, and OpenCL plus four OpenMP-target case-study directions, AST-driven augmentation for robustness testing, and a declarative spec format designed for community 
extension beyond the initial corpus. CodeRosetta~\cite{CodeRosetta2024} and HPC-Coder-V2~\cite{HPCCoderV2} demonstrate specialized modeling and fine-tuning for parallel code; ParaCodex~\cite{ParaCodex2026} introduces profiling-guided agentic translation. \parbench{} provides an execution-based substrate for evaluating such approaches against general-purpose LLMs under controlled conditions.

\paragraph{HPC-specific LLM evaluation.}
LASSI~\cite{LASSI2024} is the nearest prior work in granularity and verification, evaluating an agentic self-correction pipeline on 10 HeCBench kernels across two translation directions. \parbench{} asks a different question: what is the single-attempt translation capability across a broad set of directions? It also differs in scope---35~kernels from five suites, six standard directions among CUDA, OpenMP, and OpenCL plus four OpenMP-target case studies, AST-driven augmentation, and a declarative spec format designed for community extension. CodeRosetta~\cite{CodeRosetta2024} and HPC-Coder-V2~\cite{HPCCoderV2} demonstrate specialized fine-tuning for parallel code; ParaCodex~\cite{ParaCodex2026} introduces profiling-guided agentic translation. \parbench{} provides an execution-based substrate for evaluating all such approaches against general-purpose LLMs under controlled conditions.

% \paragraph{Benchmark validity and robustness.}
% Public code benchmarks can be memorized, and surface-form perturbations can reveal when
% evaluation depends on brittle pattern matching~\cite{SWEbenchIllusion2025}. Recent work
% formalizes this concern: CodeMorph~\cite{CodeMorph2025} applies 26~semantic-preserving
% transformations to detect data leakage in code LLM evaluation, and Chen et
% al.~\cite{MemorizeOrGeneralize2025} measure memorization risk through code rewriting on
% standard benchmarks. \parbench{}'s augmentation engine applies the same
% principle---semantics-preserving source perturbation---to the HPC domain, where
% training-data contamination is particularly likely given that suites like Rodinia appear
% in hundreds of GitHub repositories. The engine is not a complete solution because it
% cannot rule out algorithm-level familiarity, but it provides a concrete robustness probe:
% models that succeed only on unmodified source reveal surface-form dependence, while models
% that maintain pass rates under augmentation demonstrate more robust parallel reasoning.
% Its correctness-only scope also differs from efficiency-oriented work such as
% TRACY~\cite{TRACY2025}, which shows that functional correctness and performance are
% separable.

\paragraph{Benchmark validity and robustness.}
Public benchmarks can be memorized, and surface-form perturbations reveal brittle pattern matching~\cite{SWEbenchIllusion2025}. CodeMorph~\cite{CodeMorph2025} applies 26~semantic-preserving transformations to detect data leakage, and Chen et al.~\cite{MemorizeOrGeneralize2025} measure memorization risk through code rewriting. \parbench{}'s augmentation engine applies the same principle to the HPC domain, where contamination is particularly likely given that suites like Rodinia appear in hundreds of GitHub repositories. Models that succeed only on unmodified source reveal surface-form dependence; those that maintain pass rates under augmentation demonstrate more robust parallel reasoning. This correctness-only scope differs from efficiency-oriented work such as TRACY~\cite{TRACY2025}, which shows that functional correctness and runtime performance are separable objectives.

%\paragraph{Benchmark validity and robustness.}
%Public code benchmarks can be memorized, and surface-form perturbations can reveal when evaluation depends on brittle pattern matching~\cite{SWEbenchIllusion2025}. \parbench{}'s augmentation engine is not a complete solution to training-data contamination because it cannot rule out algorithm-level familiarity, but it provides a concrete robustness probe for widely circulated HPC kernels. Its correctness-only scope also differs from efficiency-oriented work such as TRACY~\cite{TRACY2025}, which shows that functional correctness and performance are separable.